\section{Setup}
\label{sec:setup}

\subsection{Research questions and the two tests}
Fine-tuning a code model on obfuscated programs raises its accuracy on those programs by roughly
nineteen points. Every question below asks what that number is made of. Adaptation can raise forward
accuracy by teaching the model to \emph{read} the program better, or by teaching it the surface a
particular transform leaves behind; only the first deserves to be called robustness. We separate them
with two tests that share no machinery, committed before any method was on the table:
\begin{itemize}[leftmargin=*, topsep=2pt, itemsep=1pt]
  \item \textbf{Composition.} Does competence on single transforms survive being stacked with a
        transform family the model has never seen?
  \item \textbf{Reversal.} Does it survive the task being run backwards --- given a return value,
        produce a call that yields it? No adapter is trained on this direction, so any backward gain
        is transferred program understanding rather than task fitting.
\end{itemize}
\textbf{RQ1} applies both to the field's existing strategies --- inference-time routing over
per-transform specialists, weight-space merging of those specialists, and broad multi-transform
fine-tuning --- treated as baselines. \textbf{RQ2} locates what adaptation latches onto instead,
because that determines the remedy. \textbf{RQ3} derives the remedy, a paired-consistency objective
anchored to the model's behaviour on the unobfuscated parent, isolates its mechanism by ablating the
teacher, and measures it across eight models. \textbf{RQ4} re-runs the two tests on the anchored
model and states the bound.

\subsection{Models}
Table~\ref{tab:models} lists the panel: eight instruct models from four lineages (Meta, BigCode,
Google, IBM), code-specialised and general, 7B to 34B, plus one pretrained checkpoint used only for
the one-shot baseline. Every adapter is a LoRA with the same rank and target modules on every model,
and one prompt builder serves training, evaluation and attention extraction.
% GENERATED 2026-09-18 00:43 UTC by scripts/analysis/59_master_tables.py
% SOURCE: configs/models.yaml
% Do not edit by hand -- regenerate. 
\begin{table*}[ht]
\centering\scriptsize\setlength{\tabcolsep}{2pt}
\caption{\textbf{The model panel.} Four lineages, code-specialised and general instruct models, 7B to 34B. \emph{Chat template}: \emph{system} accepts a system role; \emph{merged} folds the system text into the user turn; \emph{plain} has no template and is used for the one-shot baseline only. Parameter counts are nominal. Every adapter is a LoRA with the same rank and target modules on every model.}
\label{tab:models}
\begin{tabular}{@{}llllrrrl@{}}
\toprule
\textbf{Model} & \textbf{Checkpoint} & \textbf{Lineage} & \textbf{Type} & \textbf{Params (B)} & \textbf{Layers} & \textbf{Hidden} & \textbf{Chat template} \\
\midrule
CodeLlama-7B~\cite{code_llama} & \texttt{\tiny codellama/CodeLlama-7b-Instruct-hf} & Meta & coder & 6.7 & 32 & 4096 & system \\
CodeLlama-13B~\cite{code_llama} & \texttt{\tiny codellama/CodeLlama-13b-Instruct-hf} & Meta & coder & 13.0 & 40 & 5120 & system \\
CodeLlama-34B~\cite{code_llama} & \texttt{\tiny codellama/CodeLlama-34b-Instruct-hf} & Meta & coder & 33.7 & 48 & 8192 & system \\
Llama-3.1-8B~\cite{llama3herd} & \texttt{\tiny meta-llama/Llama-3.1-8B-Instruct} & Meta & instruct & 8.0 & 32 & 4096 & system \\
StarCoder2-15B~\cite{lozhkov2024starcoder2stackv2} & \texttt{\tiny bigcode/starcoder2-15b-instruct-v0.1} & BigCode & coder & 15.5 & 40 & 6144 & merged \\
Gemma-3-12B~\cite{gemma3} & \texttt{\tiny google/gemma-3-12b-it} & Google & instruct & 12.2 & 48 & 3840 & system \\
CodeGemma-7B~\cite{codegemma} & \texttt{\tiny google/codegemma-7b-it} & Google & coder & 8.5 & 28 & 3072 & merged \\
Granite-3.1-8B~\cite{granite31} & \texttt{\tiny ibm-granite/granite-3.1-8b-instruct} & IBM & instruct & 8.2 & 40 & 4096 & system \\
Llama-3.1-8B (base)~\cite{llama3herd} & \texttt{\tiny meta-llama/Llama-3.1-8B} & Meta & pretrained & 8.0 & 32 & 4096 & plain \\
\bottomrule
\end{tabular}
\end{table*}


\subsection{Dataset}
Table~\ref{tab:dataset} summarises the corpus. The split is by program, never by row; three input
cases per program are fixed and their gold outputs executed from the clean parent, so the grader is
independent of any transform. Coverage differs by condition because structural transforms decline
some programs by design and the unseen family needs enough encodable sites; every paired contrast is
computed on the programs common to the cells it compares, fixed before evaluation.
% GENERATED 2026-09-15 02:05 UTC by scripts/analysis/59_master_tables.py
% SOURCE: data/train/base/python.jsonl, data/splits/python.json, data/eval/heldout/items/*
% Do not edit by hand -- regenerate. 
\begin{table}[ht]
\centering\small
\caption{\textbf{The dataset.} One Python corpus, split by program so no obfuscated variant of a training program can reach evaluation.}
\label{tab:dataset}
\begin{tabular}{@{}p{0.34\linewidth}p{0.6\linewidth}@{}}
\toprule
\textbf{Property} & \textbf{Value} \\
\midrule
Base programs (Python) & 2231: apps 1584, cruxeval 543, humaneval 104 \\
Split & by program, seed 17: train 1563 / val 111 / held-out 557 \\
Input cases per program & 3, gold output executed from the clean parent \\
Program length & mean 8.4 LOC, max 57 \\
Held-out eval, clean & 557 programs $\times$ 3 cases = 1671 items \\
Coverage per condition & Tables~\ref{tab:taxonomy} and \ref{tab:stacks}; structural transforms decline some programs by design, the unseen family needs $\geq 3$ encodable sites \\
Training rows per adapter & one condition's train split $\times$ 3 cases (4689); breadth pools the five seen conditions \\
Paired contrasts & on the programs common to every cell in the contrast, fixed before evaluation; program-clustered bootstrap, 2{,}000 resamples \\
JavaScript & 168 held-out programs (504 items), transferred intact; not regenerable on the cluster \\
Grading & strict normalised exact match, no containment; the backward task executes the produced call against the clean parent \\
\bottomrule
\end{tabular}
\end{table}


\subsection{Obfuscation taxonomy}
Single transforms (Table~\ref{tab:taxonomy}) fall into three families: identifier (L1b, L1r, L2),
structural (S1, S2, and S2's two halves S3 and S4) and encoding (X1 and its halves). Each is applied
to the clean parent and never stacked during training. H1 is the quarantined held-out obfuscator; its
evaluation budget was spent under a pre-registered two-pass discipline, and every unseen-family number
in this paper is measured on its trainable sibling X1, which reproduces H1's difficulty at
$r = 0.9992$ (mean absolute difference 0.48 points). Stacks (Table~\ref{tab:stacks}) are a separate
namespace built by applying parts to the parent in order; order matters, and no system is trained on
a stack. Ten stacks contain only seen transforms (six at depth two, four at depth three to four); six
contain the unseen family or one of its halves and are the stimulus RQ4 turns on.
% GENERATED 2026-09-14 07:37 UTC by scripts/analysis/59_master_tables.py
% SOURCE: configs/conditions.yaml, data/eval/heldout/items/*
% Do not edit by hand -- regenerate. 
\begin{table*}[ht]
\centering\footnotesize
\caption{\textbf{Single transforms.} Each is applied to the clean parent and never stacked. \emph{Cap} is the maximum growth factor over the parent above which a program is declined rather than clipped; \emph{programs / items} is held-out coverage. H1 is the quarantined held-out obfuscator whose evaluation budget is spent; X1 is its trainable sibling and carries every unseen-family number.}
\label{tab:taxonomy}
\setlength{\tabcolsep}{3pt}\begin{tabular}{@{}llp{0.30\linewidth}lrrl@{}}
\toprule
\textbf{Code} & \textbf{Family} & \textbf{Transform} & \textbf{Trainable} & \textbf{Cap} & \textbf{Programs / items} & \textbf{Role} \\
\midrule
L0 & none & normalised original: comments and docstrings stripped, 4-space indent & yes & 1.2$\times$ & 557 / 1671 & clean reference \\
L1b & identifier & adversarial renaming of all bindings incl. the entry function, from semantic-inversion tables (max$\rightarrow$min) and an unrelated domain vocabulary & yes & 2.0$\times$ & 553 / 1659 & ladder \\
L1r & identifier & random hex renaming of all bindings (v\_a3f2, f\_9c1e) & yes & 2.0$\times$ & 557 / 1671 & ladder \\
L2 & identifier & sequential minification of all bindings (a, b, …, aa) plus type-annotation stripping & yes & 1.5$\times$ & 557 / 1671 & ladder \\
S1 & structural & control-flow flattening of the entry function into a dispatch loop with randomised state ids and shuffled cases & yes & 6.0$\times$ & 416 / 1248 & ladder \\
S2 & structural & opaque predicates (computed always-true/false guards) plus never-called dead helpers & yes & 4.0$\times$ & 556 / 1668 & ladder \\
S3 & structural & dead-code insertion only: 1--2 never-called module-level helpers (must be ignored) & yes & 4.0$\times$ & 557 / 1671 & stacks only \\
S4 & structural & opaque predicates only: 1--3 computed guards at statement boundaries (must be reasoned about) & yes & 4.0$\times$ & 556 / 1668 & stacks only \\
X1 & encoding & every string literal through an inline XOR decoder, every + $-$ \^{} through a type-guarded algebraic identity, int literals expanded & yes & 12.0$\times$ & 405 / 1215 & \textbf{unseen family} \\
X1m & encoding & X1's arithmetic half only (mechanisms: [mba]) & yes & 12.0$\times$ & 351 / 1053 & mechanism ablation \\
X1s & encoding & X1's string half only (mechanisms: [str]); site bar 1 & yes & 12.0$\times$ & 246 / 738 & mechanism ablation \\
H1 & heldout & string encoding (base64 / RC4 string array) plus guarded mixed-boolean-arithmetic rewriting & never & 12.0$\times$ & -- & quarantined \\
\bottomrule
\end{tabular}
\end{table*}

% GENERATED 2026-09-13 17:08 UTC by scripts/analysis/59_master_tables.py
% SOURCE: configs/conditions_composite.yaml, data/eval/heldout/items/*
% Do not edit by hand -- regenerate. 
\begin{table*}[ht]
\centering\footnotesize
\caption{\textbf{Stacked composites.} A separate namespace from the ladder: parts are applied to the clean parent left to right, and order matters because composition does not commute. Size caps were calibrated on a 40-program build at cap 60 (p95 $\times$ 1.25) so that the cap is never the binding constraint. No system is ever trained on a stack; the breadth and family adapters see single transforms only.}
\label{tab:stacks}
\setlength{\tabcolsep}{2.5pt}\begin{tabular}{@{}llrlrrlp{0.17\linewidth}@{}}
\toprule
\textbf{Code} & \textbf{Parts, in order} & \textbf{Depth} & \textbf{X1 inside} & \textbf{Cap} & \textbf{Prog.\,/\,items} & \textbf{Group} & \textbf{Purpose} \\
\midrule
\texttt{C\_L1b\_S1} & L1b$\to$S1 & 2 & no & 8.0$\times$ & 418 / 1254 & seen d2 & does L1b's specialist advantage survive stacking? \\
\texttt{C\_L1r\_S1} & L1r$\to$S1 & 2 & no & 8.0$\times$ & 418 / 1254 & seen d2 & headline identifier $\otimes$ structural stack \\
\texttt{C\_S1\_L1r} & S1$\to$L1r & 2 & no & 8.0$\times$ & 418 / 1254 & seen d2 & order control for C\_L1r\_S1 (renames S1's state variables) \\
\texttt{C\_L2\_S4} & L2$\to$S4 & 2 & no & 5.0$\times$ & 556 / 1668 & seen d2 & identifier $\otimes$ the reason-about-it half of S2 \\
\texttt{C\_L1r\_S3} & L1r$\to$S3 & 2 & no & 5.0$\times$ & 557 / 1671 & seen d2 & identifier $\otimes$ the ignore-it half of S2 \\
\texttt{C\_S4\_S3} & S4$\to$S3 & 2 & no & 5.0$\times$ & 556 / 1668 & seen d2 & positive control: reconstructs S2 by composition \\
\addlinespace[2pt]
\texttt{C3\_L1r\_S3\_S4} & L1r$\to$S3$\to$S4 & 3 & no & 10.0$\times$ & 556 / 1668 & seen d3--4 & depth 3, identifier first \\
\texttt{C3\_S1\_S3\_S4} & S1$\to$S3$\to$S4 & 3 & no & 14.0$\times$ & 395 / 1185 & seen d3--4 & depth 3, structural only \\
\texttt{C3\_L1r\_S1\_S4} & L1r$\to$S1$\to$S4 & 3 & no & 14.0$\times$ & 418 / 1254 & seen d3--4 & depth 3, identifier first, flattening inside \\
\texttt{C4\_L1r\_S1\_S3\_S4} & L1r$\to$S1$\to$S3$\to$S4 & 4 & no & 18.0$\times$ & 418 / 1254 & seen d3--4 & depth 4 \\
\addlinespace[2pt]
\texttt{C\_L1r\_X1} & L1r$\to$X1 & 2 & yes & 10.0$\times$ & 405 / 1215 & unseen d2 & unseen family last, after renaming \\
\texttt{C\_X1\_S1} & X1$\to$S1 & 2 & yes & 12.0$\times$ & 321 / 963 & unseen d2 & unseen material itself flattened \\
\texttt{C\_S2\_X1} & S2$\to$X1 & 2 & yes & 22.0$\times$ & 555 / 1665 & unseen d2 & unseen family applied to dead code too \\
\addlinespace[2pt]
\texttt{C\_L1r\_X1m} & L1r$\to$X1m & 2 & yes & 10.0$\times$ & 351 / 1053 & half-unseen d2 & half an unseen family (arithmetic) \\
\texttt{C\_S1\_X1s} & S1$\to$X1s & 2 & yes & 14.0$\times$ & 418 / 1254 & half-unseen d2 & half an unseen family (strings) \\
\addlinespace[2pt]
\texttt{C3\_L1r\_S1\_X1} & L1r$\to$S1$\to$X1 & 3 & yes & 16.0$\times$ & 417 / 1251 & unseen d3 & one unseen component at depth 3 \\
\bottomrule
\end{tabular}
\end{table*}


\subsection{Systems}
Every system is a LoRA on the same base model, evaluated on still-obfuscated code. \emph{Clean LoRA}
(\texttt{tuned\_L0}) is tuned on unobfuscated programs and is the control throughout; one
\emph{specialist} per seen transform; \emph{breadth} (\texttt{mono\_all}) is one adapter trained on
all five seen transforms pooled; \emph{family} (\texttt{tuned\_X1}) is trained on the unseen family's
sibling and is the ceiling for that family. \emph{Routing} mixes the specialists with a learned gate
(\texttt{mole\_router}), its argmax, a fixed uniform gate or a gate frozen at random initialisation;
the last two separate what the mixture is worth from what the routing is worth. \emph{Merging}
combines the specialists in weight space (TIES, DARE-TIES, DARE-linear) or anchors the merge on the
clean adapter. The \emph{anchored} objective (\texttt{cons\_lam3}) is Eq.~\ref{eq:paired_consistency}
with $\lambda = 3$, a plateau from a sweep on the trainable grid: the teacher is a frozen clean-code-tuned
adapter reading the unobfuscated parent, the student never sees the parent, and the objective trains on
exactly the rows breadth trains on (26{,}841 rows, 1{,}257 steps for both); at $\lambda = 0$ it reduces
to breadth identically, which makes breadth the data-matched control. A one-shot in-context baseline
completes the set.

\subsection{Measurement}
Forward accuracy is strict normalised exact match on the produced output, with no containment
matching; the backward task is graded by executing the produced call against the clean parent.
Every contrast is a paired difference on the program subset common to its cells, with a
program-clustered bootstrap of 2{,}000 resamples; bold intervals exclude zero. Decision rules were
pre-registered before submission and are reported by their verdicts even when inconclusive. A cell
whose format-failure rate exceeds 0.25 measures the prompt contract rather than the task and enters
no mean; untuned format failure measures adherence to that contract, not capability, so screening
base models by it is unsound.
